\usepackage[utf8]{inputenc}
\PassOptionsToPackage{hyphens}{url}\usepackage{hyperref}
\usepackage[landscape, left=0.5cm, right=0.5cm, top=1.2cm, bottom=0.5cm, includehead, headheight=14pt]{geometry}
\usepackage[spanish, es-tabla]{babel}
\usepackage[table,xcdraw]{xcolor}
\usepackage[italicdiff]{physics}
\usepackage{textcomp}
\usepackage{amsthm, amssymb, amsmath, amsfonts, bbm, mathtools}
\usepackage{minted}
\usepackage{enumerate}
\usepackage[shortlabels]{enumitem}
\usepackage{titling}
\usepackage{lmodern} %optimiza algunas fuentes
\usepackage{float}
\usepackage{mathrsfs}
\usepackage{wasysym}
\usepackage{csquotes}
\usepackage{dsfont}
\usepackage{caption}
\decimalpoint
\usepackage{graphicx}
\renewcommand{\baselinestretch}{1}
\usepackage{environ}

\usepackage{tocloft}
\usepackage{multicol}
\usepackage[skip=0pt]{parskip}
\usepackage{fancyhdr}
\usepackage{titlesec}

%%%%%%%%%%%%%%%%%% tabla de contenidos %%%%%%%%%%%%%%%%%%
\setlength{\cftbeforesecskip}{1pt} % Espacio entre secciones en el índice
\renewcommand{\cftsecdotsep}{\cftdotsep} % Poner puntitos a las secciones principales
\renewcommand{\cftsecfont}{\normalfont\footnotesize} % Letra más chica
\renewcommand{\cftsecpagefont}{\normalfont\footnotesize}
\renewcommand{\cftsubsecfont}{\normalfont\scriptsize}
\renewcommand{\cftsubsecpagefont}{\normalfont\scriptsize}

\newcommand{\Pa}[1]{\left( #1 \right)}
\newcommand{\PA}[1]{\big( #1 \big)}

%Images path
\graphicspath{{img/}}

%%%%%%%%%%%%%%%%%% d EXTENDED VERSION %%%%%%%%%%%%%%%%%%
\usepackage{etoolbox}
\newif\ifextendedversion
\extendedversiontrue   % Descomenta para la versión EXTENDIDA (Final Nacional)
% \extendedversionfalse  % Descomenta para la versión COMPACTA (Hiperregional)

\newcommand{\ext}[1]{\ifextendedversion#1\fi}

% 3. Bloque largo / Entorno (múltiples párrafos, listas, problemas tipo)
\NewEnviron{extended}{
    \ifextendedversion
        \vspace{0.3em}
        \begingroup
        \small\color{black!80}
        \vspace{0.2em}
        \noindent\textit{\footnotesize Versión Extendida:} \vspace{-1mm}\par
        \BODY
        \vspace{0.2em}
        \endgroup
        \vspace{0.3em}
    \fi
}


%%%%%%%%%%%%%%%%%% Minted %%%%%%%%%%%%%%%%%%
% \fontsize{tamaño de letra}{espacio entre líneas}\selectfont
\newcommand{\codesize}{\fontsize{7pt}{8pt}\selectfont} % tamfuente e interlineado para 9pt en el doc con extarticle
\setminted{
    fontsize=\codesize,      % Usa \scriptsize o \tiny si tienes muy buena vista
    baselinestretch=0.85,      % Comprime el espacio vertical entre líneas
    linenos=true,              % Números de línea (esenciales para debuggear en papel)
    numbersep=3pt,             % Pega los números de línea al código
    xleftmargin=12pt,          % Reduce el margen izquierdo que minted agrega por defecto
    breaklines=true,           % Ajuste de línea automático si te pasas del ancho
    breakanywhere=true,        % Rompe la línea incluso si no hay espacios
    tabsize=2,                 % Indentación corta (vital en 3 columnas)
    style=vs,                  % Tema visual claro y de alto contraste al imprimir
    stripall=true              % Elimina líneas en blanco al inicio y al final del archivo
}

% \inputcode{lenguaje}{ruta}{nombre_a_mostrar}
\newcommand{\inputcode}[3]{
    {\raggedleft \scriptsize \ttfamily \detokenize{#3} \par} 
    \vspace{-4pt}
    \inputminted[firstline=1, firstnumber=1]{#1}{#2}
}

%%%%%%%%%%%%%%%%%% Explicacion de codigo %%%%%%%%%%%%%%%%%%

% \algoinfo{Tiempo}{Memoria}{Explicación rápida}
\newcommand{\algoinfo}[3]{
    % \vspace{-2pt}
    \vspace{-1pt}
    {\fontsize{8pt}{9pt}\selectfont \textbf{T:} $O(#1)$ $\mid$ \textbf{M:} $O(#2)$ $\mid$ \textit{#3} \par}
    % \vspace{2pt}
}

%%%%%%%%%%%%%%%%%% Numeraciones %%%%%%%%%%%%%%%%%%
% Figuras numeradas
\usepackage{chngcntr}
\counterwithin{figure}{section}

% Ecuaciones numeradas
\numberwithin{equation}{section}

\allowdisplaybreaks

%%%%%%%%%%%%%%%%%% Definiciones de mate %%%%%%%%%%%%%%%%%%
\renewcommand{\theenumi}{\alph{enumi}}
\renewcommand{\labelenumi}{{\theenumi})}
\renewcommand\qedsymbol{$\blacksquare$}
\newcommand{\N}{\mathbb{N}}
\newcommand{\Z}{\mathbb{Z}}
\newcommand{\R}{\mathbb{R}}
\newcommand{\I}{\mathbb{I}}
\newcommand{\Q}{\mathbb{Q}}
\newcommand{\Zm}[1]{\mathbb{Z}/ #1 \mathbb{Z}}
\newcommand\scalemath[2]{\scalebox{#1}{\mbox{\ensuremath{\displaystyle #2}}}}

%%%%%%%%%%%%%%%%%% Comandos especiales de mate %%%%%%%%%%%%%%%%%%
\newcommand{\pa}[1]{\left( #1 \right)}
\newcommand{\br}[1]{\left[ #1 \right]}
\DeclarePairedDelimiter{\set}{\{}{\}}
\DeclarePairedDelimiter{\floor}\lfloor\rfloor
\DeclarePairedDelimiter{\ceil}\lceil\rceil
\DeclareMathOperator{\sgn}{sgn}
\DeclareMathOperator{\ord}{ord}
\newenvironment{solution}
{\renewcommand\qedsymbol{$\square$}\begin{proof}[Solución]}
    {\end{proof}}

%%%%%%%%%%%%%%%%%% Espaciado %%%%%%%%%%%%%%%%%%
\setlength{\columnsep}{0.25cm} % Espacio minimo entre columnas
\setlength{\columnseprule}{0.5pt} % Linea divisoria opcional para guiar la vista
\setlength{\parindent}{0pt}
\setlist[enumerate,1]{label=\arabic*., nosep}
\setlist[itemize]{noitemsep}
% \newcommand{\noenter}[1]{\vspace{-2mm}}


% \titlespacing*{comando}{izq}{arriba}{abajo}
% Espaciado hiper-compacto
\titlespacing*{\section}{0pt}{5pt}{3pt}
\titlespacing*{\subsection}{0pt}{3pt}{1pt}
\titlespacing*{\subsubsection}{0pt}{2pt}{1pt}

% Jerarquía visual (Pensada para impresión monocromática de ICPC)

% SECTION: MAYÚSCULAS y una línea negra sólida por debajo
\titleformat{\section}{\normalfont\normalsize\bfseries}{}{0pt}{\MakeUppercase}[{\vspace{1pt}\titlerule[0.8pt]}]

% SUBSECTION: Flecha indicadora, negrita (sin línea)
\titleformat{\subsection}{\normalfont\normalsize\bfseries}{}{0pt}{$\blacktriangleright$\hspace{3pt}}

% SUBSUBSECTION: Cursiva negrita y punto discreto (letra un poco más chica)
\titleformat{\subsubsection}{\normalfont\small\bfseries\itshape}{}{0pt}{$\bullet$\hspace{3pt}}

% Eliminar el salto de línea post-titulo (Run-in titles) si quieres compresión extrema:
% \titlespacing*{\subsection}{0pt}{1pt}{1em} % El tercer valor en 'em' hace que el texto siga en la misma línea

% Header
\pagestyle{fancy}
\headheight 14pt % Reducido de 25pt
\fancyhead{}
\fancyfoot{}
\lhead{\textbf{Reference - ICPC} $\mid$ Mäthgic Crüe $\mid$ UG-CIMAT $\mid$ Junio 2026} % Todo en una línea
\rhead{\ifextendedversion { \large \textbf{-- Versión Extendida} } \fi \textbf{Pág. \thepage}}
\renewcommand{\headrulewidth}{0.4pt}
\setlength\headsep{4pt} % Pegar el texto principal al header

% Comprimir espacio vertical alrededor de ecuaciones
\AtBeginDocument{
  \setlength{\abovedisplayskip}{2pt}
  \setlength{\belowdisplayskip}{2pt}
  \setlength{\abovedisplayshortskip}{0pt}
  \setlength{\belowdisplayshortskip}{0pt}
}